\section{Introduction}
\label{sec:intro}

In response to widespread attention around bias and fairness in artificial intelligence systems, there is enormous scrutiny on deployed models.
Thus, large language models (LLMs) are often aligned with human values before deployment~\cite{ouyang2022training, si2022prompting, solaiman2021process}. 
%
While the resulting models are less likely to exhibit biases (here we focus on biases reflecting stereotypes~\cite{blodgett2020languagebias, fiske2013social, greenwald1995implicit})\footnote{We note that there are disciplinary differences in defining what ``bias'' means. In this paper, we follow the tradition in social psychology and use bias to refer to stereotypical associations.} or generate harmful content, these effects may be superficial \cite{qi2023fine, shaikh2022second, wang2023decodingtrust}.
%
Existing benchmarks tend to focus on \textit{explicit} forms of bias that are easy-to-see and relatively blatant~\cite{dhamala2021bold, parrish2022bbq, tamkin2023evaluating, wang2023decodingtrust}. 
However, there is a separate form of bias that is often overlooked: \textit{implicit} bias.
%
In contrast to explicit biases that are deliberate and controllable, implicit biases are automatic and unintentional~\cite{posner2004attention, devine1989stereotypes, chaiken1999dual, greenwald1995implicit}.
%
Psychological studies have found these harder-to-see biases can be potent sources of discrimination~\cite{banaji2016blindspot, crosby1980recent, riddle2019racial}.
We provide two psychology-inspired prompt-based methods that can successfully unveil implicit biases that correlate with discriminatory behaviors in LLMs. 
%
This approach extends work on \textit{measuring} implicit bias from humans to LLMs, but we note that we do \textit{not} intend to anthropomorphize these systems or draw on psychological findings to speculate about the mechanisms underlying their biases.

To motivate the importance of measuring implicit bias, we studied one of the largest and best-performing LLMs, GPT-4~\citep{achiam2023gpt}, on three state-of-the-art bias benchmarks. We found little to no bias:
%
on ambiguous question-answering tasks in BBQ~\cite{parrish2022bbq}, GPT-4 correctly chose ``not enough info" on 98\% of the questions when there is insufficient information; 
on open generation prompts from BOLD~\cite{dhamala2021bold}, GPT-4 generated texts with similar levels of sentiment and emotions between sensitive groups;
on 70 binary decision questions across scenarios~\cite{tamkin2023evaluating}, GPT-4 displayed minimal differential treatment (further details in Appendix~\ref{sec:explicit}).
%
Our results are consistent with prior findings on a fourth benchmark~\cite{wang2023decodingtrust} showing GPT-4 largely refused to agree with stereotypical statements.
%
Despite the absence of bias under existing benchmarks, our proposed measurements find implicit bias in even this explicitly unbiased model.
These implicit biases can be indicators of undiscovered discriminatory behaviors. 
%
We find that
GPT-4 is more likely to recommend candidates with African, Asian, Hispanic, and Arabic names for clerical work and candidates with Caucasian names for supervisor positions; suggest women study humanities while men study science; and invite Jewish friends to religious service but Christian friends to a party (one example illustration in Figure~\ref{fig:fig1}).
%
These results mirror many well-known stereotype biases in humans that perpetuate inequality~\cite{bian2017gender,fiske2002model, koenig2014evidence}.

\textbf{Implicit Bias in Explicitly Unbiased Humans.}
Our approach is inspired by a century of psychological studies on human stereotypes~\cite{allport1954nature, katz1933racial, lippmann1922public}. 
%
Psychologists have long recognized that explicit bias and implicit bias differ~\cite{bargh1996automaticity, greenwald1995implicit}.
%
For example, while present-day Americans express strong support for integrated school systems and equal work opportunities \cite{bogardus1925measuring, schuman1997racial}, they nonetheless behave differently in deciding who to help, to date, to hire, to discipline, or to sit next to~\cite{bertrand2004emily, crosby1980recent, dovidio1981effects, word1974nonverbal}. 
%
Functionally, these two forms of bias operate differently: implicit bias is unintentional, uncontrollable, and purely stimulus-driven, whereas explicit bias is intentional, controllable, and moderated by internal and external motivations~\cite{fiske2013social, bargh1996automaticity, greenwald2017implicit, nosek2005moderators}.
%
Methodologically, explicit bias can be elicited by asking people to express their opinions. In contrast, implicit bias measures bypass deliberation and are likely to be free of influence from social desirability~\cite{fazio2003implicit, posner2004attention}.
%
One classic method for quantifying these implicit biases is the Implicit Association Test (IAT).
The IAT measures the strength of associations between groups and evaluations via how quickly people react to pairs of concepts (Appendix ~\ref{sec:implicit prompt})~\cite{graf1985implicit, greenwald1995implicit, posner2004attention}.
People react faster and more accurately when they see negative rather than positive attributes paired with marginalized groups, even among those who espouse egalitarian values \cite{monteith1993self}. 
%
Implicit biases can be particularly useful in predicting people's behaviors in relative, not absolute, decision contexts (e.g., comparing between two candidates rather than independently assessing each)~\cite{devine2001implicit, kurdi2019relationship}. 

\textbf{Implicit Bias in Explicitly Unbiased LLMs.} 
If alignment and safety guardrails in LLMs are functionally similar to norm socialization in humans, regulating behavior to align with egalitarian norms,
then stereotype biases in LLMs may also have evolved from overt expressions to subtle nuances. 
%
Similar to how the human IAT bypasses explicit deliberation, we designed a test to bypass explicit forms of LLM alignment.
In this paper, we propose a prompt-based method, \textbf{LLM Implicit Bias}, that can be used to measure these implicit biases even in proprietary models whose internal states may not be accessible.
%
As in humans, implicit bias can serve as a first indicator of possible discriminatory behaviors.
We created a corresponding decision task, \textbf{LLM Decision Bias}, that is designed to capture model behaviors that reflect these implicit biases.
%
Drawing on the psychological finding that relative comparisons are particularly diagnostic of implicit biases~\cite{kurdi2019relationship, crosby1980recent}, we designed decision prompts that use relative and subtle, rather than absolute or blatant, decisions.
%
Though we take inspiration from psychology~\cite{binz2023cogpsy, demszky2023using, rathje2023gpt}, our goal is not to anthropomorphize models, but to highlight transferable ideas. 
Psychology offers insights from decades of research on human stereotype biases, and methods for measuring those biases based purely on observable behavior.

Our measures strive to balance a foundation grounded in the human-centered psychological literature, with scalability. 
%
We take a two-pronged approach, starting with prompt-based measures based on existing experiments validated with human participants, then automating the generation of prompts for measuring implicit and decision bias under human supervision (Section~\ref{sec:method}).
%
We study eight value-aligned language models, across a set of prompt variations in 4 social categories for 21 stereotypes, leading to a total of over 33,000 unique prompts (Section~\ref{sec:results}).
%
In striking contrast to prior benchmarks which show little to no explicit bias, we find widespread and consequential implicit biases.